\section{Alerting}

Alerts are a part of observability that interrupts people. That makes them different from
dashboards---a bad dashboard is inconvenient, but a bad alert wakes someone up at 3\,am for the
wrong reason, or worse, stays silent when something actually needs attention.

The platform stores alert rules as code, following the same review and promotion process as
everything else \cite{lsst-it-k8s-cookbook,prometheus-alerting,prometheus-alertmanager}. But the
technical side is the easier part. The harder part is writing alerts that are actually useful.

A good alert is not just a PromQL expression that fires when a number crosses a threshold
\cite{prometheus-alerting,prometheus-alertmanager}. It carries enough context to start a diagnosis:
what the condition is, how urgent it is, who should receive it, and where to look first. Without
that, an alert becomes noise---something that fires, gets acknowledged, and gets ignored. With it,
an alert is the first step of a structured response.

There is also a timing problem specific to night operations. During an observing shift, not
everything that is wrong needs immediate human attention. A saturated disk, a failed exporter, a
flapping network link, and a degraded storage component are serious, but they are not all the same
kind of serious. The alerting system needs to reflect that, routing the right signal to the right
person at the right time \cite{prometheus-alertmanager}, without flooding the on-call operator with
everything at once. Getting that routing right is ongoing work, and it improves the same way
everything else does: through real incidents, honest postmortems, and small adjustments that
accumulate over time.

% \textbf{TODO: ADD ALERTING STATISTICS HERE}